\chapter{Conclusion and Future Scope}
\label{ch:conclusion}

\section{Conclusion}

The project set out to answer a single operational question: when an
LLM--driven AutoML system presents a natural--language or
semi--structured plan and then emits a Python pipeline, does the
emitted code mechanically implement the plan it claims to realise?
Existing evaluations check that the code compiles, that it runs,
that it does not raise exceptions and that it reaches an acceptable
score on some metric. None of these checks answers the alignment
question, which is strictly stronger and which decides whether the
documentation, approvals and deployment gates that users see
actually describe the pipeline that is running.

To address this, the project has:

\begin{itemize}
    \item defined \emph{reasoning--code alignment} as a stand--alone
    verification problem that is cleanly decoupled from reasoning
    correctness, program correctness, execution success and
    predictive performance;
    \item designed a Reasoning Intermediate Representation of
    atomic typed claims, a Program Intermediate Representation that
    fuses abstract syntax, dataflow and library--level semantics,
    and a mechanical mapping between the two;
    \item specified coverage, weighted alignment, critical mismatch
    and localisation measures, together with an AutoML--specific
    mismatch taxonomy covering missing, contradictory,
    model--level, parameter--level, ordering, evaluation,
    data--handling, constraint, unsupported, dead--reasoning and
    extraneous categories;
    \item proposed an evaluation protocol against execution--only,
    performance--only, LLM--judge and keyword/API baselines on
    natural LLM outputs and on seeded--mutation corpora with
    human--annotated targets, with a statistical design that
    reports false--alignment rates, macro--F1 of category detection,
    and localisation accuracy with appropriate confidence
    intervals; and
    \item engineered a running end--to--end implementation --- five
    FastAPI services and a Streamlit front--end, a Redis Queue worker, a
    configurable LLM runner (Ollama in the reported runs), a bounded
    subprocess executor, a first faithfulness verifier, a PostgreSQL
    database with
    local--filesystem artefact storage and RQ1--RQ5 analysis handlers;
    and
    \item evaluated that implementation on a real 270--cell sweep,
    where the naive and schema prompt conditions were perfectly
    reliable, the meta--feature condition was markedly less reliable
    and more expensive without improving predictive score, and the
    mechanical verifier flagged seven pipelines that executed and
    scored well but whose reasoning did not match their code.
\end{itemize}

The objectives stated in Chapter~\ref{ch:intro} are therefore met at
three distinct levels, which the report keeps separate.
\emph{Design and formalisation} of the full reasoning--code alignment
framework is complete. \emph{Software} exists and runs for the
pipeline and a first faithfulness verifier (action presence and
meta--feature value grounding) on \texttt{main}. \emph{Empirical
evidence} exists for the prompting study and for the core hypothesis
that execution success does not imply alignment (H1), with preliminary
support for discrimination (H2). What remains is not design but
construction and measurement: the graph--based verifier, independent
pre--code reasoning capture, and the seeded--mutation corpus needed to
quantify the baseline comparison (H3, H4 and the full RQ2/RQ5
metrics). The project thus delivers a validated foundation and a
precise, bounded agenda for the remaining work, rather than a
completed verifier presented as finished.

Taken together, these contributions position the proposed
verification layer as a complement, not a replacement, for existing
AutoML evaluation. Where execution evidence tells us that a
pipeline \emph{can} succeed, reasoning--code alignment tells us
whether it has succeeded for the \emph{reason} it was supposed to.

\section{Limitations}

The present system has several limitations that bound the strength
of its claims.

\textbf{Plan expressiveness.} Natural--language plans may be vague,
incomplete or even impossible to execute. Structured claims decrease
ambiguity, but they also change the reasoning distribution and
ignore helpful context. A learnt RIR extractor may itself reflect
the extractor's interpretation of the plan rather than the LLM's
intent.

\textbf{Static analysis of Python.} The arbitrary behaviour of
Python cannot be determined by static analysis alone, because of
reflection, dynamic imports, metaprogramming, services, callbacks
and notebook state limiting graph recovery. Runtime traces only
exercise paths that were taken. General semantic equivalence is
undecidable; bounded evidence via behavioural probes necessarily
leaves such cases as \textsc{unverifiable}.

\textbf{Ontology and ranking.} The proposed verifier relies on
heuristic operator--ontology mapping and candidate ranking, which lean
towards common libraries and may not cover custom code well. High
alignment is also not a guarantee of quality: a pipeline can be
faithful to its plan yet the plan itself be poor, the data biased or
the model poorly calibrated. Finally, the emitted reasoning does not
reveal the model's internal computation; the project checks
artefact--to--artefact conformance, not mechanistic faithfulness.

\textbf{Post--hoc reasoning trace.} The B2 reasoning trace is printed
by the generated code at run time, so it is a post--hoc self--report
rather than a plan captured before the code. The verifier can catch a
trace that contradicts its own code, but not a case where the code and
the trace are wrong in the same way (for example, both naming
\texttt{StandardScaler} when the dataset called for
\texttt{RobustScaler}). Independent pre--code reasoning capture
(\texttt{B2R}) is designed but not implemented, and is the single most
important next step.

\textbf{Scope.} The first version of the system targets
compositional scikit--learn tabular pipelines. Deep--learning,
distributed training, multi--modal workflows and notebooks with
hidden state remain outside scope.

\textbf{LLM supply--chain.} Hosted models, prompts, APIs and
libraries change over time; alignment numbers must be interpreted
alongside the model identifiers, dates, temperatures, seeds and
package versions recorded at run time.

\section{Future Scope}

The project opens several promising directions.

\textbf{Deeper program semantics.} Symbolic execution and SMT
constraints for parameters and branches, blending static checks with
runtime monitors and causal traces, can push the verifier past the
current static--analysis bound.

\textbf{Typed pipeline IRs.} Building typed pipeline IRs for
proof--carrying AutoML would turn the current best--effort witness
mapping into a certified contract between the planner and the coder.

\textbf{Broader domains.} Extensions comprise deep--learning and
multimodal pipelines, contracts between agents, deployment--time
checks and certified generation.

\textbf{Repair loops.} An optimising repair loop can take the
localised mismatch and ask a coder to patch only the diverged
subgraph. This must respect the rule set fixed during evaluation;
high--impact \textsc{unverifiable} claims should still be subject
to human approval.

\textbf{Integration into AutoML products.} The verifier can be
offered as a drop--in library inside existing AutoML products so
that any pipeline that is deployed can be audited claim by claim
against the plan that was approved. Such mechanisms would not, by
definition, make AutoML correct, but would at least allow for one
promise to be audited: that the code executed implements the
pipeline claimed by the system.

\section{Final Remarks}

Execution tests can miss a verification problem with LLM--driven
AutoML. A system may proclaim a consistent strategy for
preprocessing, model choice, ordering or evaluation, produce code that
runs and gives a passable score, and still not implement the plan it
presented; reasoning--code alignment cannot be replaced by execution
correctness and predictive performance. This project formalised
alignment, proposed the reasoning and program intermediate
representations, defined the mechanical verdicts, metrics, taxonomy
and baselines, and --- crucially --- built a running system and
measured it. On that first 270--cell run, seven pipelines that
executed and scored well were caught contradicting their own
reasoning: direct evidence that the divergence is real and
mechanically detectable. The open empirical question is now narrower
and sharper --- how \emph{often} successful AutoML code differs from
its plan across models, tasks and complexity, and how precisely the
divergence can be localised --- and the software, protocol and first
measurements are in place to answer it.
